\chapter{From P1 to the Dual DC-SCA Subproblem}

\section{P2: covariance lifting and SDR}
Introduce $\mat{W}_k=\vect{w}_k\vect{w}_k^H$. The equality is equivalent to
$\mat{W}_k\succeq0$ plus $\operatorname{rank}(\mat{W}_k)=1$. Dropping the rank
constraint yields an SDP relaxation. This is a relaxation, not an assertion
that rank one is automatically guaranteed in the present coupled robust ISAC
problem.

\section{Robust and PCRB LMIs}
The S-procedure converts the norm-bounded worst-case communication constraint
to an LMI with a nonnegative multiplier. The trace-PCRB constraint is expressed
through a Schur block LMI using an auxiliary matrix $\mat{M}_p$ and
$\operatorname{tr}(\mat{M}_p)\leq\Gamma_{\rm track,p}$. The full state dimension
$N_\theta$ is retained in this block; reducing it to an unrelated scalar block
would change the model.

\section{P3: double DC-SCA}
The rank penalty is based on the gap between trace and dominant eigenvalue; the
binary penalty is based on $b_{mp}(1-b_{mp})$. At each iteration, concave terms
are linearized around the previous iterate, producing a convex SDP subproblem.
The approach is a stationary-point heuristic with penalty continuation, not a
global certificate for the original mixed-integer problem.

\begin{figure}[H]
\centering
\includegraphics[width=0.82\linewidth]{fig5_statistical_double_dc_convergence.png}
\caption{Continuous double-DC stabilization over 25 common realizations.
The figure supports support stabilization before discrete recovery; it is not
itself a physical-feasibility result.}
\end{figure}

\section{Why the recovery step is separate}
The stable Top-$N_{\rm req}$ support of a continuous relaxation suggests a
binary topology but does not prove it physically feasible. Therefore, the
algorithm projects to binary $b$, solves a fixed-assignment continuous problem,
extracts rank-one beamformers when needed, and runs an explicit validator. This
separation is scientifically important: it prevents fractional relaxation
quality from being reported as a deployable result.

\sourcebox{Derivation source: \texttt{CONVEXIFICATION\_CHAIN.md}; caveats:
\texttt{SDR\_TIGHTNESS.md} and \texttt{COMPLEXITY\_BOUND.md}.}
